%%%%%%%%%%  scan idx 56  |  folio 43  %%%%%%%%%%
tel que
pour tout $E\in\cC_0$ les $F^i(E)$ soient dans $\cC'_0$ et nuls sauf un nombre fini, on en déduit un homomorphisme de $\underline{K}(\cC_0)$ dans $\underline{K}(\cC'_0)$. Ces remarques s'étendent aux multifoncteurs~; de plus, les suites spectrales «~classiques~» donnent des propriétés de transitivité que nous n'expliciterons pas.

En particulier, si $F$ est un foncteur exact, la formule (2.11) se simplifie :
\begin{equation}\tag{2.11 bis}
 F(\gamma_{\cC_0}(E))=\gamma_{\cC'_0}(F(E)).
\end{equation}
Ceci s'applique en particulier si $\cC=\cC'$, si $F$ est le foncteur identique et si $\cC_0\subset\cC'_0$~: on a un homomorphisme canonique de $\underline{K}(\cC_0)$ dans $\underline{K}(\cC'_0)$.

\textbf{Théorème 2.3.} \emph{Soient $\cC$ une catégorie abélienne et $\cC_0\subset\cC_1$ deux sous-classes telles que les classes} (à isomorphisme près) \emph{d'objets de $\cC_1$ forment un ensemble. On suppose les conditions suivantes satisfaites}
\begin{enumerate}[label=(\roman*)]
\item \emph{Pour toute suite exacte $0\to A'\to A\to A''\to0$ dans $\cC$ telle que $A$ et $A''$ soient dans $\cC_0$ (resp. $\cC_1$), on a $A'\in\cC_0$ (resp. $A'\in\cC_1$).}
\item \emph{Soient $u:A\to B$ et $v:C\to B$ avec $A,B,C\in\cC_1$ et $u$ surjectif~; alors il existe $L\in\cC_0$ et des homomorphismes $u':L\to C$ et $v':L\to A$ tels que $u'$ soit surjectif et $vu'=uv'$.}
\item \emph{Pour tout $A\in\cC_1$, il existe un entier $d=d(A)$ tel que pour toute résolution gauche $L$ de $A$ par des objets de $\cC_0$, on ait $Z_d(\underline{L})\in\cC_0$} ($Z_d$ désigne les cycles de dimension $d$).
\end{enumerate}
\emph{Dans ces conditions, l'homomorphisme canonique de $K(\cC_0)$ dans $K(\cC_1)$ est un isomorphisme.}

\textbf{Principe de la démonstration~:} soit $A\in\cC_1$~; il existe en vertu de (ii) et (iii) une résolution finie $\underline{L}$ de $A$ par des objets de $\cC_0$~; posons
